% DOCUMENT SETUP
\documentclass[letterpaper,11pt]{article}
\hyphenpenalty=8000
\textwidth=125mm
\textheight=200mm
\usepackage[top=3cm, bottom=3cm, inner=3cm, outer=3cm, includehead]{geometry}
\usepackage{fancyhdr}
\setlength{\headheight}{13.6pt} % Fix for fancyhdr warning
\addtolength{\topmargin}{-1.6pt} % Optional: further fix as suggested
\pagestyle{fancy}
\fancyhead{}
\fancyfoot{}
\raggedbottom
\usepackage{xurl}
\usepackage{graphicx}
\usepackage{xcolor}
\usepackage{amssymb}
\usepackage{mathtools}
\usepackage{alltt}
\usepackage{amsmath}
\usepackage[pdftex, hidelinks]{hyperref}
\urlstyle{same}
\usepackage[T1]{fontenc}
\usepackage[utf8]{inputenc}
\usepackage{lmodern}
\usepackage{csquotes}
\usepackage[notes, backend=biber]{biblatex-chicago}
\pagenumbering{arabic}
\setcounter{page}{1}
\graphicspath{{./imgs/}}
\usepackage[english]{babel}
\usepackage{multirow}
\usepackage{tabularx}
\usepackage{fancyvrb}
\usepackage{caption}
%\usepackage[french]{babel}
%\usepackage[spanish]{babel}

% AUTHOR AND TITLE
\begin{document}
\fancyhead[RO]{ELEC 475 Lab 2, Alistair Barfoot and Luke Barry\ \ \ \ \thepage}
\begin{center}
  \LARGE
  \textbf{ELEC 475 Lab 2, Alistair Barfoot and Luke Barry}\\[12pt]
  \normalsize
\end{center}
\tableofcontents

\newpage

\section{Network Architecture}
% You should describe the structure in sufficient detail so that it can be reimplemented solely from the textual description. You should also have proper citations to the techniques that you used.

\begin{table}[h]
  \centering
  % \resizebox{\linewidth}{!}{
  \begin{tabularx}{\linewidth}{|c|c|c|X|}
    \hline
    \textbf{Layer}    & \textbf{Input Shape} & \textbf{Output Size} & \textbf{Parameters}                                                                                      \\
    \hline
    \textbf{Conv1}    & 227x227x3            & 114x114x64           & \verb+in_channels=3+, \verb+out_channels=64+, \verb+kernel_size=3+, \verb+stride=2+, \verb+padding=1+    \\
    \hline
    \textbf{ReLU1}    & 114x114x64           & 114x114x64           &                                                                                                          \\
    \hline
    \textbf{MaxPool1} & 114x114x64           & 57x57x64             & \verb+kernel_size=3+, \verb+stride=2+, \verb+padding=1+                                                  \\
    \hline
    \textbf{Conv2}    & 57x57x64             & 29x29x128            & \verb+in_channels=64+, \verb+out_channels=128+, \verb+kernel_size=3+, \verb+stride=2+, \verb+padding=1+  \\
    \hline
    \textbf{ReLU2}    & 29x29x128            & 29x29x128            &                                                                                                          \\
    \hline
    \textbf{MaxPool2} & 29x29x128            & 15x15x128            & \verb+kernel_size=3+, \verb+stride=2+, \verb+padding=1+                                                  \\
    \hline
    \textbf{Conv3}    & 15x15x128            & 8x8x256              & \verb+in_channels=128+, \verb+out_channels=256+, \verb+kernel_size=3+, \verb+stride=2+, \verb+padding=1+ \\
    \hline
    \textbf{ReLU3}    & 8x8x256              & 8x8x256              &                                                                                                          \\
    \hline
    \textbf{MaxPool3} & 8x8x256              & 4x4x256              & \verb+kernel_size=3+, \verb+stride=2+, \verb+padding=1+                                                  \\
    \hline
    \textbf{FC1}      & 4x4x256(4096)        & 1024                 &                                                                                                          \\
    \hline
    \textbf{ReLU4}    & 1024                 & 1024                 &                                                                                                          \\
    \hline
    \textbf{FC2}      & 1024                 & 1024                 &                                                                                                          \\
    \hline
    \textbf{ReLU5}    & 1024                 & 1024                 &                                                                                                          \\
    \hline
    \textbf{FC3}      & 1024                 & 2 (u,v)              &                                                                                                          \\
    \hline
  \end{tabularx}
  % }
\end{table}


\end{document}